\section{\methodname{}: Justifying the Practical Implementation}
\label{app:matnormal}
Here we detail how, starting from~\Cref{sec:sample-loss-lmo}, we arrive at the practical  implementation of \methodname{} in \Cref{alg:ours}.  


\subsection{Estimating the Gradients}
\label{sec:weighted-estimation}
Let $G_s \in \mathbb{R}^{d_{\mathrm{out}} \times d_{\mathrm{in}}}$ be a batch gradient estimate at iteration $s$. This is the gradient with respect to the full parameter matrix $W_s \in  \mathbb{R}^{d_{\mathrm{out}} \times d_{\mathrm{in}}}.$ 
Let $G_{A,s}$ and $G_{B_s}$ be the mini-batch factor gradients at iteration $s$. 
Though in~\eqref{e-grad-wrt-factors}  we denoted the factor gradients as the full batch estimates, in reality we only have access to stochastic mini-batch estimates. To estimate the full batch factor gradients we use an EMA (Exponential Moving Average) defined as 
\begin{align}
    M_A &= \mathrm{EMA}_{s\leq t}^{\beta_1}[G_{A,s}] \\
    M_B &= \mathrm{EMA}_{s\leq t}^{\beta_1}[G_{B,s}] \label{eq:mom-grad-factors}
\end{align}
where for any sequence $d_s$ we define the exponential moving average as:
\[\mathrm{EMA}_{s\leq t}^{\beta}[d_s] : = \sum_{s \leq t }(1-\beta_2)\beta_2^{t-s} d_s. \]
This definition is such it is equivalent to computing $ M_A $ and $ M_B$ as done in~\Cref{alg:ours}.
% We estimate both weighted moments of \Cref{sec:weighted-lmo}
% from the factor gradients alone, taking the momentum $\widehat M_t$ for
% the objective and the model factors $P^\star,Q^\star$ for the
% constraint.
% , with each replacement of a past-step factor by its current
% value made under \Cref{as:slow-drift}.


\subsection{Estimating the Preconditioners}
\label{sec:kronecker-factor-estimates}
There are two huddles to estimating the $P$ and $Q$ factors in \methodname{} as needed in the update~\eqref{eq:polorafull-update} and detailed in~\Cref{sec:sample-loss-lmo}. First, we do not have direct access to $\Sigma$ as 
defined in~\Cref{lem:gradprop}. Second, we need to approximate $\Sigma$ as a product $Q  \otimes P$~\eqref{eq:pq-factor} using only access to the factor gradients $G_A$ and $G_B.$ Here we detail how we circumvent these issues, and how exactly we arrive at the $p,q$
recursion of \Cref{alg:ours} and recovering the
update~\eqref{eq:polora-update}. We will also detail in~\Cref{sec:scale-placement} 
why \Cref{alg:ours} normalizes the accumulated vectors $p,q$ after the
average.
% The
% update is invariant to the scale of the curvature factors, but the
% average that estimates them depends on that scale.


First, the factors $P$ and $Q$ are a Kronecker product approximation of $\Sigma$~\eqref{eq:pq-factor}, but we do not have access to $\Sigma$. This is because $\Sigma$ is an average over the outerproduct of gradients each evaluated over one input (see~\Cref{lem:gradprop}). These single input gradients are expensive to store, and backpropagtion gives only the average these gradients over a batch. To circumvent this issue, instead of computing an average over individual gradients, we approximate $\Sigma$ as an exponential moving average over  batch gradients. That is, let $g_s = \vec(G_s)$ be the batch gradient at iteration $s$, and let
\begin{equation} \label{eq:SigmaEMA}
    \Sigma_t = \mathrm{EMA}_{s\leq t}^{\beta_2}[ g_s g_s^\top ]
\end{equation}
We never actually compute this $\Sigma_t$ (it would be impossible to store), but rather, we will use~\eqref{eq:SigmaEMA} from now on as the new definition of $\Sigma_t$, from which we will compute the preconditioners $Q_t$ and $P_t.$


To estimate the Kronecker factorization~\eqref{eq:pq-factor}, we first assume both $P = \diag(p)$ and $Q=diag(q)$.  Thus our objective now is to find $p_t$, and $q_t$  such that
  \begin{equation}\label{e-diag-kron-second-mom-model}
    \mathrm{EMA}_{s\leq t}^{\beta_2}[g_sg_s^\top]\approx (\diag(q_t) \otimes \diag(p_t)).
  \end{equation}

Our main issue now, is that in LoRA training, we do not have access to $g_s = \vec(G_s)$. Instead, backpropagation gives the factor gradients $G_{A,s}=B_s^\top G_s$
and $G_{B,s}=G_sA_s^\top$~\eqref{e-grad-wrt-factors} at no extra cost. Forming
the full gradient $G_s$, by contrast, would take a $d_{\mathrm{out}}\times
d_{\mathrm{in}}$ matrix multiplies and full-matrix storage per layer, exactly what
the low-rank structure avoids. Consequently we need to estimate the preconditioners $Q$ and $P$ using only the factor gradients.



To estimate $P$ and $Q$ from the factor gradients alone, we link their second moments to the second moment of the batch gradient given in~\eqref{eq:SigmaEMA}. 
To this end, note that
$\vec(G_{A,s})=(I\otimes B_s^\top)\vec(G_s)$. Consequently the second moment of the gradient of $A$ is given by
\begin{align*}
 \Sigma_{A,t} &:=  \mathrm{EMA}_{s \leq t}^{\beta_2}\big[ \vec(G_{A,s})  \vec(G_{A,s})^\top\big]  \\
 &=  \mathrm{EMA}_{s \leq t}^{\beta_2}\big[ (I\otimes B_s^\top)\vec(G_s)\vec(G_s)^\top   (I\otimes B_s)\big]  \\
   &\approx  (I\otimes B_t^\top) \mathrm{EMA}_{s \leq t}^{\beta_2}\big[ g_s g_s^\top   \big]  (I\otimes B_t),
\end{align*}
where the approximation replaces $B_s$ by the current $B_t$, so that it can move outside of the EMA.
 Using~\eqref{e-diag-kron-second-mom-model} in the above gives
\begin{align} \label{eq:SigmaAapprox}
 \Sigma_{A,t} 
   &\approx    (I\otimes B_t^\top)  (Q_t\otimes P_t)  (I\otimes B_t) \; =     Q_t\otimes  (B_t^\top P_t B_t),
\end{align} 
where the equality uses the mixed-product rule of Kronecker products.
Using analogous arguments for the gradient of $B$, we arrive at the following approximations
\begin{align}
  \Sigma_{A,t} & = \mathrm{EMA}_{s \leq t}^{\beta_2}\big[\vec(G_{A,s})\vec(G_{A,s})^\top\big] \approx Q_t\otimes C_{A,t} \nonumber \\
  \Sigma_{B,t}&=\mathrm{EMA}_{s \leq t}^{\beta_2}\big[\vec(G_{B,s}^\top)\vec(G_{B,s}^\top)^\top\big] \approx  P_t\otimes C_{B,t}, \label{eq:observable-moments}
\end{align}
where $C_{A,t}:=B_t^\top P_tB_t$ and $C_{B,t}:=A_tQ_tA_t^\top$ are the $r\times r$ matrices of~\eqref{eq:ours-closure} evaluated at the model factors.

\paragraph{Fitting the factors.} We will now use~\eqref{eq:observable-moments} as our evidence for fitting $P$ and $Q$.
% For brevity, here we drop the subscript $t$ since
% all quantities are at the current step.
Following
KL-Shampoo~\cite{klshampoo}, we fit with the log-determinant
divergence~\cite{logdetdiv}
\[
  D(\Sigma,S)=\Tr(\Sigma S^{-1})-\log\det(\Sigma S^{-1})-d,
\]
defined for positive definite matrices $\Sigma$ and $S$ of common
dimension $d$. Let $\widehat{P}$ and $\widehat {Q}$ be our current estimates of $P$ and $Q$.
% Away from the exact model, no candidate reproduces the
% observed moment, and the divergence decides how the mismatch is
% weighed. 
Using~\eqref{eq:observable-moments} we fit the two factors by solving
\begin{equation}\label{eq:factor-fit}
  \hat{q}(\widehat P)
  =\argmin_{q>0}\;
  D\big(\Sigma_{A,t},\; \diag(q) \otimes \widehat C_A\big),
  \qquad
  \hat{p}(\widehat Q)
  =\argmin_{p>0}\;
  D\big(\Sigma_{B,t},\;\diag(p) \otimes \widehat C_B\big),
\end{equation}
where we freeze the current estimate of $\widehat C_A$ when estimating $\hat{q}$, and freeze the estimate of $\widehat C_B$ when estimating $\hat{p}$. This alternating freezing is needed to get a closed form update for $\hat{q}$ and $\hat{p}$, since otherwise $\widehat C_A$ and $\widehat C_B$ introduce a joint dependency through
$\widehat C_A=B^\top\widehat PB$ and
$\widehat C_B=A\widehat QA^\top$, for diagonal
$\widehat P,\widehat Q\succ0$.


The closed form solutions to~\eqref{eq:factor-fit} are given by
\begin{equation}
    \label{eq:factor-fit-sol}
      \hat{q}(\widehat P)=\frac1r\diag\!\left(\mathrm{EMA}_{s\leq t}^{\beta_2}[G_{A,s}^\top \widehat C_A^{-1}G_{A,s}]\right),
  \qquad
  \hat{p}(\widehat Q)=\frac1r\diag\!\left(\mathrm{EMA}_{s\leq t}^{\beta_2}[G_{B,s} \widehat C_B^{-1}G_{B,s}^\top]\right).
\end{equation}

For fixed positive definite $\Sigma$ and $S$, let
\begin{equation} \label{eq:opt-scalar-div}
c^\star(\Sigma,S) := \tfrac{1}{d}\Tr(\Sigma S^{-1}) =\argmin_{c>0}D(\Sigma,cS)\end{equation}
be the best scalar multiple
of $S$.
The
next lemma shows the coupled fits are consistent even when posed with
inexact estimates.
\nikhil{I think this lemma is not useful. In fact it suggests you do not need a stateful estimator of $p$, $q$, you could use $\hat{q}(1)$ and $\hat{p}(1)$ for example.}
\begin{lemma}[Consistency]\label{lem:consistency}
Assume that~\eqref{e-diag-kron-second-mom-model}  holds with equality. Furthermore, assume that there exists $p^\star, q^\star, C_A$ and $C_B$ such that
\begin{align}
  \Sigma_{A,t} & =   \diag(q^\star) \otimes C_{A} \nonumber \\
  \Sigma_{B,t}&=  \diag(p^\star) \otimes C_{B}, \label{eq:observable-moments-lem}
\end{align}
For any
diagonal $\widehat P\succ0$ and $\widehat Q\succ0$, the
fits~\eqref{eq:factor-fit}  with $\widehat C_A=B^\top\widehat PB$
and $\widehat C_B=A\widehat QA^\top$ recover the true $p^\star$ and $q^\star$ upto a multiplicative constant:
\begin{equation} \label{eq:awdawda}
  \hat{q}(\widehat P)= c^\star\big(C_A,\widehat C_A\big)\,\diag(q^\star),
  \qquad
  \hat{p}(\widehat Q)= c^\star\big(C_B,\widehat C_B\big)\,\diag(p^\star).
\end{equation}
\end{lemma}
\begin{proof}
Since we assume~\eqref{eq:observable-moments} holds to equality, we have that $\Sigma_{A,t}$ is block
diagonal with $j$th block $ q_j^\star\,C_A$, and the candidate
$\diag(q)\otimes\widehat C_A$ of~\eqref{eq:factor-fit} at a point $q>0$
is block diagonal with $j$th block $q_j\,\widehat C_A$. Traces and
log-determinants add across blocks, so the objective is
$\sum_jD(  q^\star_jC_A,\,q_j\widehat C_A)$, and thus is separable 
over the unknowns $q_j$. Using~\eqref{eq:opt-scalar-div} solving for the scalar $q_j$ gives the best scalar
multiple:
\[
\hat{q}_j=c^\star( q^\star_jC_A,\widehat C_A)
=\,q^\star_j\,c^\star(C_A,\widehat C_A).
\]
The factor
$\,c^\star(C_A,\widehat C_A)$ is the same for every $j$, and consequently~\eqref{eq:awdawda} holds. The claim
for $\hat{p}$ follows with the roles of the two factors exchanged.
\end{proof}


\Cref{lem:consistency} shows that, if the second moments are block diagonal~\eqref{eq:observable-moments-lem}, which given the approximation in~\eqref{eq:SigmaAapprox} is reasonable, then by fitting our estimates of $P$ and $Q$ using~\eqref{eq:factor-fit-sol}, we will exactly recover the true $P^\star$ and $Q^\star$ up to a multiplicative constant. Since we know our estimates can be off by a multiplicative constant, we simply re-normalize them so that $\|q\|_{\infty} = \|p\|_{\infty} =1.$ This avoids the factors drifting across iterations. Putting all of the argument above together, we arrive at our online estimates for $P$ and $Q$ as follows.

\nikhil{the "online estimator" still needs explanation / justification?}
\rob{What part still needs justifying? The EMA is well justified, since it is the solution in~\eqref{eq:factor-fit-sol}}
\nikhil{not true. the EMA uses $C_{B,s}$ from $s < t$, whereas as \eqref{eq:factor-fit-sol} uses $C_{B,t}$.}

\paragraph{Online estimator.}

% \rob{** OLd ***}
%  Since the product 
%   $Q \otimes P $ is unchanged by the rescaling
%   $(P,Q)\mapsto(aP,a^{-1}Q)$, we fix
%   $\norm{\diag P}_\infty=\norm{\diag Q }_\infty=1$ 
Unrolling the EMAs in~\eqref{eq:factor-fit-sol}, with the post normalization given by \Cref{lem:consistency} gives
\begin{equation}\label{eq:ours-klcoupling}
\begin{aligned}
  \hat{q}_t(P_t) &= \frac1r\diag\!\left(G_{A,t}^\top C_{A,t}^{-1}G_{A,t}\right),
  &
  \hat{p}_t(Q_t) &= \frac1r\diag\!\left(G_{B,t} C_{B,t}^{-1}G_{B,t}^\top\right),
  \\
  q_{t+1} &= \beta_2q_t+(1-\beta_2)\hat{q}_t(P_t),
  &
  p_{t+1} &= \beta_2p_t+(1-\beta_2)\hat{p}_t(Q_t),
  \\
  Q_{t+1} &= \diag\!\bigl(q_{t+1}/\norm{q_{t+1}}_\infty\bigr),
  &
  P_{t+1} &= \diag\!\bigl(p_{t+1}/\norm{p_{t+1}}_\infty\bigr).
\end{aligned}
\end{equation}

%  Unrolled, the recursion weights the per-step estimates by
% exactly the $w_{t,s}$ of \Cref{sec:weighted-lmo}. Within the window,
% $B_s\approx B_t$ by \Cref{as:slow-drift}, and
% the stored $P_s$ changes only at rate $1-\beta_2$, so
% $C_{A,s}\approx C_{A,t}$ and each per-step estimate already uses the
% current matrices. The running $q_t$ then matches the full fit
% $\hat{q}(P_t)$, the consistent fit of \Cref{lem:consistency}.



\paragraph{Recovering the update of \Cref{alg:ours}.} 
Putting everything together, replacing the full batch factor gradients in
the update~\eqref{eq:factor-directions} with their momentum estimates~\eqref{eq:mom-grad-factors} with one additional "look ahead" step look-ahead step gives line~\ref{ln:Bspectral} in \Cref{alg:ours}. Using~\eqref{eq:ours-klcoupling} to compute our approximation for $P$ and $Q$ gives lines~2,~\ref{ln:qmom} and~\ref{ln:pmom}  
 The only departure of \Cref{alg:ours}
from this derivation is that the preconditioners update after the
direction step, so step $t$ runs on the estimate from step $t-1$, which
avoids a second formation and inversion of $C_A$ and $C_B$ per step.

\subsection{Normalization Placement}
\label{sec:scale-placement}
\methodname{} normalizes the accumulated vectors $q_{t+1},p_{t+1}$ to unit
largest entry when it forms the next-step factors
$Q_{t+1},P_{t+1}$~\eqref{eq:ours-klcoupling}, rather than normalizing each
per-step estimate $\hat{q}_t,\hat{p}_t$ before it enters the average.
Normalizing the accumulated vectors does not change the update, which
depends only on the direction of the stored factors. Normalizing each
estimate first would change the average, which is sensitive to the scale
of those estimates.

\rob{This repeats statement}
\paragraph{The update is scale-free.} The update depends on the stored
factors only through their direction. Under a rescaling $Q\mapsto cQ$
($c>0$), both $Q^{-1/2}$ and $C_B^{-1/2}=(AQA^\top)^{-1/2}$ pick up a
factor $c^{-1/2}$. Inside the matrix sign this scalar has no effect, since
the sign is scale-invariant. Outside it, the scalar multiplies $D_A$ and
$D_B$ by $c^{-1/2}$, which the division by $\norm{D_A}_2$ and
$\norm{D_B}_2$ in the magnitude step cancels. The factor $P$ enters the
same way. The update therefore sees only the direction of the stored
factors, never the magnitude, so normalizing the accumulated
vectors leaves it unchanged.

\paragraph{The average is scale-sensitive.} The coupled fit makes each
per-step estimate scale inversely with the stored factor it is computed
from,
\[
  \hat{q}_t(aP_t)=a^{-1}\,\hat{q}_t(P_t),
  \qquad
  \hat{p}_t(aQ_t)=a^{-1}\,\hat{p}_t(Q_t)
  \qquad(a>0),
\]
since $C_A=B^\top PB$~\eqref{eq:ours-closure} is linear in $P$, and
likewise $C_B$ in $Q$. Normalizing after
averaging gives every stored factor the same scale
($\norm{\diag P_t}_\infty=1$ at every step), so the estimates enter the
average on a common footing. Were the stored factor at step $s$ left carrying a step-varying scale
$a_s$ instead, the identity above would scale
$\hat{q}_s$ by $1/a_s$, so it would enter with effective weight
$w_{t,s}/a_s$ rather than the $w_{t,s}$ chosen in \Cref{sec:weighted-lmo}.
Normalizing the per-step estimates themselves before averaging would fail
for a second reason. Normalization is nonlinear and does not commute with
the mean, so per-step noise that cancels in the raw average would survive.
This average-then-normalize order is the one by which momentum precedes
normalization in normalized SGD~\cite{cutkosky_nigt} and orthogonalization
in Muon~\cite{denoise_ortho}.

